\documentclass[a4paper,12pt,twoside]{memoir}

% Castellano
\usepackage[spanish,es-tabla]{babel}
\selectlanguage{spanish}
\usepackage[utf8]{inputenc}
\usepackage[T1]{fontenc}
\usepackage{lmodern} % Scalable font
\usepackage{microtype}
\usepackage{placeins}

\RequirePackage{booktabs}
\RequirePackage[table]{xcolor}
\RequirePackage{xtab}
\RequirePackage{multirow}

% Links
\usepackage[colorlinks]{hyperref}
\hypersetup{
	allcolors = {red}
}

% Ecuaciones
\usepackage{amsmath}

% Rutas de fichero / paquete
\newcommand{\ruta}[1]{{\sffamily #1}}

% Párrafos
\nonzeroparskip


% Imagenes
\usepackage{graphicx}
\newcommand{\imagen}[2]{
	\begin{figure}[!h]
		\centering
		\includegraphics[width=0.9\textwidth]{#1}
		\caption{#2}\label{fig:#1}
	\end{figure}
	\FloatBarrier
}

\newcommand{\imagenflotante}[2]{
	\begin{figure}%[!h]
		\centering
		\includegraphics[width=0.9\textwidth]{#1}
		\caption{#2}\label{fig:#1}
	\end{figure}
}



% El comando \figura nos permite insertar figuras comodamente, y utilizando
% siempre el mismo formato. Los parametros son:
% 1 -> Porcentaje del ancho de página que ocupará la figura (de 0 a 1)
% 2 --> Fichero de la imagen
% 3 --> Texto a pie de imagen
% 4 --> Etiqueta (label) para referencias
% 5 --> Opciones que queramos pasarle al \includegraphics
% 6 --> Opciones de posicionamiento a pasarle a \begin{figure}
\newcommand{\figuraConPosicion}[6]{%
  \setlength{\anchoFloat}{#1\textwidth}%
  \addtolength{\anchoFloat}{-4\fboxsep}%
  \setlength{\anchoFigura}{\anchoFloat}%
  \begin{figure}[#6]
    \begin{center}%
      \Ovalbox{%
        \begin{minipage}{\anchoFloat}%
          \begin{center}%
            \includegraphics[width=\anchoFigura,#5]{#2}%
            \caption{#3}%
            \label{#4}%
          \end{center}%
        \end{minipage}
      }%
    \end{center}%
  \end{figure}%
}

%
% Comando para incluir imágenes en formato apaisado (sin marco).
\newcommand{\figuraApaisadaSinMarco}[5]{%
  \begin{figure}%
    \begin{center}%
    \includegraphics[angle=90,height=#1\textheight,#5]{#2}%
    \caption{#3}%
    \label{#4}%
    \end{center}%
  \end{figure}%
}
% Para las tablas
\newcommand{\otoprule}{\midrule [\heavyrulewidth]}
%
% Nuevo comando para tablas pequeñas (menos de una página).
\newcommand{\tablaSmall}[5]{%
 \begin{table}
  \begin{center}
   \rowcolors {2}{gray!35}{}
   \begin{tabular}{#2}
    \toprule
    #4
    \otoprule
    #5
    \bottomrule
   \end{tabular}
   \caption{#1}
   \label{tabla:#3}
  \end{center}
 \end{table}
}

%
% Nuevo comando para tablas pequeñas (menos de una página).
\newcommand{\tablaSmallSinColores}[5]{%
 \begin{table}[H]
  \begin{center}
   \begin{tabular}{#2}
    \toprule
    #4
    \otoprule
    #5
    \bottomrule
   \end{tabular}
   \caption{#1}
   \label{tabla:#3}
  \end{center}
 \end{table}
}

\newcommand{\tablaApaisadaSmall}[5]{%
\begin{landscape}
  \begin{table}
   \begin{center}
    \rowcolors {2}{gray!35}{}
    \begin{tabular}{#2}
     \toprule
     #4
     \otoprule
     #5
     \bottomrule
    \end{tabular}
    \caption{#1}
    \label{tabla:#3}
   \end{center}
  \end{table}
\end{landscape}
}

%
% Nuevo comando para tablas grandes con cabecera y filas alternas coloreadas en gris.
\newcommand{\tabla}[6]{%
  \begin{center}
    \tablefirsthead{
      \toprule
      #5
      \otoprule
    }
    \tablehead{
      \multicolumn{#3}{l}{\small\sl continúa desde la página anterior}\\
      \toprule
      #5
      \otoprule
    }
    \tabletail{
      \hline
      \multicolumn{#3}{r}{\small\sl continúa en la página siguiente}\\
    }
    \tablelasttail{
      \hline
    }
    \bottomcaption{#1}
    \rowcolors {2}{gray!35}{}
    \begin{xtabular}{#2}
      #6
      \bottomrule
    \end{xtabular}
    \label{tabla:#4}
  \end{center}
}

%
% Nuevo comando para tablas grandes con cabecera.
\newcommand{\tablaSinColores}[6]{%
  \begin{center}
    \tablefirsthead{
      \toprule
      #5
      \otoprule
    }
    \tablehead{
      \multicolumn{#3}{l}{\small\sl continúa desde la página anterior}\\
      \toprule
      #5
      \otoprule
    }
    \tabletail{
      \hline
      \multicolumn{#3}{r}{\small\sl continúa en la página siguiente}\\
    }
    \tablelasttail{
      \hline
    }
    \bottomcaption{#1}
    \begin{xtabular}{#2}
      #6
      \bottomrule
    \end{xtabular}
    \label{tabla:#4}
  \end{center}
}

%
% Nuevo comando para tablas grandes sin cabecera.
\newcommand{\tablaSinCabecera}[5]{%
  \begin{center}
    \tablefirsthead{
      \toprule
    }
    \tablehead{
      \multicolumn{#3}{l}{\small\sl continúa desde la página anterior}\\
      \hline
    }
    \tabletail{
      \hline
      \multicolumn{#3}{r}{\small\sl continúa en la página siguiente}\\
    }
    \tablelasttail{
      \hline
    }
    \bottomcaption{#1}
  \begin{xtabular}{#2}
    #5
   \bottomrule
  \end{xtabular}
  \label{tabla:#4}
  \end{center}
}



\definecolor{cgoLight}{HTML}{EEEEEE}
\definecolor{cgoExtralight}{HTML}{FFFFFF}

%
% Nuevo comando para tablas grandes sin cabecera.
\newcommand{\tablaSinCabeceraConBandas}[5]{%
  \begin{center}
    \tablefirsthead{
      \toprule
    }
    \tablehead{
      \multicolumn{#3}{l}{\small\sl continúa desde la página anterior}\\
      \hline
    }
    \tabletail{
      \hline
      \multicolumn{#3}{r}{\small\sl continúa en la página siguiente}\\
    }
    \tablelasttail{
      \hline
    }
    \bottomcaption{#1}
    \rowcolors[]{1}{cgoExtralight}{cgoLight}

  \begin{xtabular}{#2}
    #5
   \bottomrule
  \end{xtabular}
  \label{tabla:#4}
  \end{center}
}


















\graphicspath{ {./img/} }

% Capítulos
\chapterstyle{bianchi}
\newcommand{\capitulo}[2]{
	\setcounter{chapter}{#1}
	\setcounter{section}{0}
	\chapter*{#2}
	\addcontentsline{toc}{chapter}{#2}
	\markboth{#2}{#2}
}

% Apéndices
\renewcommand{\appendixname}{Apéndice}
\renewcommand*\cftappendixname{\appendixname}

\newcommand{\apendice}[1]{
	%\renewcommand{\thechapter}{A}
	\chapter{#1}
}

\renewcommand*\cftappendixname{\appendixname\ }

% Formato de portada
\makeatletter
\usepackage{xcolor}
\newcommand{\tutor}[1]{\def\@tutor{#1}}
\newcommand{\course}[1]{\def\@course{#1}}
\definecolor{cpardoBox}{HTML}{E6E6FF}
\def\maketitle{
  \null
  \thispagestyle{empty}
  % Cabecera ----------------
\noindent\includegraphics[width=\textwidth]{cabecera}\vspace{1cm}%
  \vfill
  % Título proyecto y escudo informática ----------------
  \colorbox{cpardoBox}{%
    \begin{minipage}{.8\textwidth}
      \vspace{.5cm}\Large
      \begin{center}
      \textbf{TFG del grado en Ingeniería Informática}\vspace{.6cm}\\
      \textbf{\LARGE\@title{}}
      \end{center}
      \vspace{.2cm}
    \end{minipage}

  }%
  \hfill\begin{minipage}{.20\textwidth}
    \includegraphics[width=\textwidth]{escudoInfor}
  \end{minipage}
  \vfill
  % Datos de alumno, curso y tutores ------------------
  \begin{center}%
  {%
    \noindent\LARGE
    Presentado por \@author{}\\ 
    en Universidad de Burgos --- \@date{}\\
    Tutores: \@tutor{}\\
  }%
  \end{center}%
  \null
  \cleardoublepage
  }
\makeatother

\newcommand{\nombre}{Marta Monje Blanco} %%% cambio de comando

% Datos de portada
\title{Sistema de etiquetado de Fitolitos Online}
\author{\nombre}
\tutor{Dr. Alvar Arnaiz González y Dr. José Francisco Díez Pastor}
\date{\today}

\begin{document}

\maketitle


\newpage\null\thispagestyle{empty}\newpage


%%%%%%%%%%%%%%%%%%%%%%%%%%%%%%%%%%%%%%%%%%%%%%%%%%%%%%%%%%%%%%%%%%%%%%%%%%%%%%%%%%%%%%%%
\thispagestyle{empty}


\noindent\includegraphics[width=\textwidth]{cabecera}\vspace{1cm}

\noindent Dr. Alvar Arnaiz González y Dr. José Francisco Díez Pastor, profesores del departamento de Ingeniería Civil, área de Lenguajes y Sistemas Informáticos.

\noindent Exponen:

\noindent Que la alumna D.ª \nombre, con DNI 71312200X, ha realizado el Trabajo final de Grado en Ingeniería Informática titulado título de TFG. 

\noindent Y que dicho trabajo ha sido realizado por el alumno bajo la dirección del que suscribe, en virtud de lo cual se autoriza su presentación y defensa.

\begin{center} %\large
En Burgos, {\large \today}
\end{center}

\vfill\vfill\vfill

% Author and supervisor
\begin{minipage}{0.45\textwidth}
\begin{flushleft} %\large
Vº. Bº. del Tutor:\\[2cm]
Dr. Alvar Arnaiz González
\end{flushleft}
\end{minipage}
\hfill
\begin{minipage}{0.45\textwidth}
\begin{flushleft} %\large
Vº. Bº. del co-tutor:\\[2cm]
Dr. José Francisco Díez Pastor
\end{flushleft}
\end{minipage}
\hfill

\vfill

% para casos con solo un tutor comentar lo anterior
% y descomentar lo siguiente
%Vº. Bº. del Tutor:\\[2cm]
%D. nombre tutor


\newpage\null\thispagestyle{empty}\newpage




\frontmatter

% Abstract en castellano
\renewcommand*\abstractname{Resumen}
\begin{abstract}
La arqueobotánica es la sub-rama de la arqueología que estudia la relación entre plantas y humanos en el pasado mediante el análisis de restos vegetales hayados en un contexto arqueológico. Su objetivo es de dar a conocer qué técnicas de agricultura y de cosecha se usaban, en qué condiciones paleoambientales subsistían y la finalidad para las que las cultivaban (fines médicos, alimenticios...) en función de la antigüedad y el estado de las muestras.

En este proyecto se va a trabajar con la recogida de datos y en el prototipado de modelos para la futura detección y clasificación de micro-restos, más concretamente de fitolitos. Los fitolitos son células de origen vegetal las cuales han sufrido un proceso de biomineralización debido a la actividad metabólica.

La finalidad de este proyecto es la de proporcionar un conjunto de herramientas y de prototipos que permitan la futura integración de un sistema de clasificación automática de fitolitos, clasificándolos en siete clases diferentes.

Con el fin de facilitar la recogida de datos para el entrenamiento del modelo se ha utilizado un sistema de cliente servidor en el cual se proporciona un etiquetador especializado con el cual los expertos podrán subir sus propias imágenes y eiquetarlas. Una vez se tuviera un conjunto razonablemente grande se emplearía para entrenar el clasificador.

Actualmente dicho servicio está desplegado en un servidor para ofrecer una mejor comunicación entre los expertos en arqueobotánica y la gestión del entrenamiento del clasificador. 
\end{abstract}

\renewcommand*\abstractname{Descriptores}
\begin{abstract}
Servidor web, Reconocimiento y clasificación de imágenes, Arqueobotánica, Fitolito, Inteligencia artificial, tratamiento y recogida de datos.
\end{abstract}

\clearpage

% Abstract en inglés
\renewcommand*\abstractname{Abstract}
\begin{abstract}Archaeobotany is the sub-branch of archeology that studies the relation between plants and humans in the past by analyzing plant residues in an archaelogical context. Its objective is to reveal what agriculture and harvesting techniques were used, in what paleoenvironmental conditions subsisted and how they were used based on the age of the samples.
	

In this project we will work with data collection and prototyping of models for the future detection and classification of micro-residues, more specifically phytoliths. Phytoliths are cells of plant origin which have undergone a process of biomineralization due to metabolic activity.

The objective of this project is to approach the classification of different photographic samples of micro-residues taken in the laboratory, trying to recognize phytoliths and classifying them into seven different classes.

In order to facilitate the collection of training data of the model, a client-server system has been used in which a specialized labelling tool is provided with which experts can upload their own images and label them. Once a reasonably large set was available, it would be used to train the classifier.

Currently, this service is deployed on a server to offer better communication among experts in archaeobotany and the management of classifier training.
\end{abstract}

\renewcommand*\abstractname{Keywords}
\begin{abstract}

Web server, Images recognition and classification, Archaeobotany, Phytolite, Artificial intelligence, treatment and data collection.
\end{abstract}

\clearpage

% Indices
\tableofcontents

\clearpage

\listoffigures

\clearpage

\listoftables
\clearpage

\mainmatter
\capitulo{1}{Introducción}
Los fitolitos son un tipo de células vegetales que se encuentran en estado de fosilización debido a los procesos metabólicos que realiza. Este estado de fosilización permite que se pueda extraer mucha información de ellas.

Actualmete existe un número muy limitado de expertos que sepan reconocer con claridad los distintos fitolitos que pueden aparecer en una muestra.

Los fitolitos son un células tridimensionales, que al ser tratadas como imágenes de dos dimensiones pueden variar la forma en la que los percibimos. Esta ambigüedad hace que el coste de tener a un ser humano etiquetando no sea beneficioso.

A día de hoy es fácil encontrar sistemas de detección y clasificación en muchos ámbitos laborales.

Hay ocasiones en las que se necesita precisión y los sistemas de detección automáticos. Actualmente se utilizan este tipo de sistemas para diagnóstico médico, detección de alimentos en mal estado, piezas defectuosas en una cadena de producción.

En este proyecto se abordan dos frentes diferenciados: La recogida e interpretación de datos y el despliegue de esta herramienta en un servidor de producción. 
\begin{itemize}
	\item \textbf{Etiquetador:} se trata de una herramienta especializada mediante la cual se puden etiquetar siete tipos predefinidos de fitolitos. Su finalidad es la de facilitar la recogida de datos de entrenamiento, haciéndola más transparente para los expertos que van a general el conjunto de datos de entrenamiento y más accesible para recoger dichos datos y manipularlos posteriormente.

	\item \textbf{Despliegue:} Las herramientas que se creen en este proyecto serán alojadas en un servidor de producción\footnote{Enlace de acceso: \url{http://martamonjeblanco.pythonanywhere.com/}}. 
	De esta forma hacemos que la aplicación sea más accesible para los usuarios	evitando instalaciones intermedias además de facilitar el paso de la información desde el etiquetador a la posterior interpretación y uso de los datos.
	
\end{itemize}

Parte del objetivo de este proyecto era la de crear un clasificador que detectase los fitolitos de forma automática y que los clasificase según su tipo, aunque por diversas razones, explicadas en el apartado de \ref{c.5}, no ha sido posible.

 La idea principal de este proyecto empieza en el etiquetado. El etiquetador genera conjuntos de etiquetas seguros gracias a la limitación en la nomenclatura de las etiquetas. Estas etiquetas, junto a las imágenes que son usadas para el etiquetado, son recogidos del servidor y se usan para generar la información necesaria para el entrenamiento de un modelo.


\capitulo{2}{Objetivos del proyecto}

Este apartado explica de forma precisa y concisa cuales son los objetivos que se persiguen con la realización del proyecto. Se puede distinguir entre los objetivos marcados por los requisitos del software a construir y los objetivos de carácter técnico que plantea a la hora de llevar a la práctica el proyecto.

\capitulo{3}{Conceptos teóricos}

En aquellos proyectos que necesiten para su comprensión y desarrollo de unos conceptos teóricos de una determinada materia o de un determinado dominio de conocimiento, debe existir un apartado que sintetice dichos conceptos.

Algunos conceptos teóricos de \LaTeX \footnote{Créditos a los proyectos de Álvaro López Cantero: Configurador de Presupuestos y Roberto Izquierdo Amo: PLQuiz}.

\section{Secciones}

Las secciones se incluyen con el comando section.

\subsection{Subsecciones}

Además de secciones tenemos subsecciones.

\subsubsection{Subsubsecciones}

Y subsecciones. 


\section{Referencias}

Las referencias se incluyen en el texto usando cite \cite{wiki:latex}. Para citar webs, artículos o libros \cite{koza92}.


\section{Imágenes}

Se pueden incluir imágenes con los comandos standard de \LaTeX, pero esta plantilla dispone de comandos propios como por ejemplo el siguiente:

\imagen{escudoInfor}{Autómata para una expresión vacía}



\section{Listas de items}

Existen tres posibilidades:

\begin{itemize}
	\item primer item.
	\item segundo item.
\end{itemize}

\begin{enumerate}
	\item primer item.
	\item segundo item.
\end{enumerate}

\begin{description}
	\item[Primer item] más información sobre el primer item.
	\item[Segundo item] más información sobre el segundo item.
\end{description}
	
\begin{itemize}
\item 
\end{itemize}

\section{Tablas}

Igualmente se pueden usar los comandos específicos de \LaTeX o bien usar alguno de los comandos de la plantilla.

\tablaSmall{Herramientas y tecnologías utilizadas en cada parte del proyecto}{l c c c c}{herramientasportipodeuso}
{ \multicolumn{1}{l}{Herramientas} & App AngularJS & API REST & BD & Memoria \\}{ 
HTML5 & X & & &\\
CSS3 & X & & &\\
BOOTSTRAP & X & & &\\
JavaScript & X & & &\\
AngularJS & X & & &\\
Bower & X & & &\\
PHP & & X & &\\
Karma + Jasmine & X & & &\\
Slim framework & & X & &\\
Idiorm & & X & &\\
Composer & & X & &\\
JSON & X & X & &\\
PhpStorm & X & X & &\\
MySQL & & & X &\\
PhpMyAdmin & & & X &\\
Git + BitBucket & X & X & X & X\\
Mik\TeX{} & & & & X\\
\TeX{}Maker & & & & X\\
Astah & & & & X\\
Balsamiq Mockups & X & & &\\
VersionOne & X & X & X & X\\
} 

\capitulo{4}{Técnicas y herramientas}
\label{c.4}

A lo largo del desarrollo del proyecto se han hecho uso de técnicas y de herramientas que han ayudado en las tareas de desarrollo y diseño de la aplicación.

\section{Etiquetador}
Para poder generar el conjunto de entrenamiento necesitamos un etiquetador que nos proporcione una salida interpretable. Preferiblemente se quería que el formato de salida fuese legible para los humanos. 

Al principio del proyecto se pensó en encontrar un etiquetador que estuviese a parte del proyecto, es decir, que fuese una aplicación distinta y que el usuario que quisiera etiquetar imágenes usase dicha aplicación, para que más adelante pudiésemos entrenar el modelo y que este fuese lo único que hubiese en el servidor.

Más tarde se optó por embeber un etiquetador dentro de la aplicación para facilitar al usuario el etiquetado y evitar posibles problemas con la posterior lectura. 

Previamente a la elección del etiquetador que se ba a usar se exploraron tras opciones:

\begin{itemize}
<<<<<<< HEAD
	\item \textbf{LabelID:}\footnote{Enlace a la páina de LabelID: \url{https://sweppner.github.io/labeld/}}:Es una aplicación de escritorio. En un principio parecía atractiva, pero debido a la poca documentación que tenía y a la forma en la que recogía los datos, se descartó.
=======
	\item \textbf{LabelID}\footnote{Enlace a la páina de LabelID: \url{https://sweppner.github.io/labeld/}}:Es una aplicación de escritorio. En un principio parecía atractiva, pero debido a la poca documentación que tenía y a la forma en la que recogía los datos, se descartó.
>>>>>>> master

	\item \textbf{LabelImg:} \footnote{Enlace a la página de LabelImg:  \url{https://github.com/tzutalin/labelImg}} Personalmente, a pesar de no haber escogido este, es probablemente el que mejor devuelva los datos. Si bien es cierto que es difícil de integrar en el proyecto, Hay mucha documentación a cerca de como entrenar modelos de Tensorflow con este etiquetador, lo cual es un punto a favor porque no tienes que manipular los datos de salida demasiado, asi que es más difícil corromperlos.
	
	\item \textbf{VGG Image Annotator (VIA):}\footnote{Enlace para acceder a VGG Image Annotator (VIA)\url{http://www.robots.ox.ac.uk/\%7Evgg/software/via/}} este etiquetador es de software libre, por lo que es fácil de integrar y de modificar dentro del proyecto. Devuelve los datos en \textit{.csv} y en \textit{.json} y la recogida de estos se puede manipular para amoldarla a la que necesite el modelo para ser entrenado.
	
	Al final fue la herramienta que se escogió y en base a la que se ha construido el etiquetador (Ver imagen \ref{fig:etiqueta}). El resultado final del etiquetador el funcional y seguro para que no se manipulen los daos de forma externa.
	
	\imagen{etiqueta}{Vista del etiquetador una vez a sido integrado en el proyecto.} 
\end{itemize}

\section{Despliegue}
\subsection{Framework : \textit{Flask}}
Para desarrollar una aplicación web necesitamos antes de nada un framework que nos provea de las utilidades propias y un esquema de trabajo. Existen diferentes alternativas, en este caso desarrolladas sobre \textit{Python}, como podrían ser \textit{Flask}, \textit{Pyramid} o \textit{Django}. Fuera de contexto ninguna es mejor que otra, pero desde el punto de vista de la aplicación y el mío propio es adecuado hacer un pequeño análisis de elección.

La decisión está entre\textit{ Django} y \textit{Flask}, por la documentación y los tutoriales existentes, no obstante, aunque \textit{Django} sea más completo, también es menos flexible. Por otro lado \textit{Flask} te permite instalar diferentes utilidades a medida que las vayas necesitando, además de que aunque tiene una documentación más reducida que \textit{Django}, esto la hace más práctica y útil dado que es lo suficientemente completa.

 \textit{Flask} también es más rápida que \textit{Django} y más sencilla de aprender, por lo tanto emplearemos \textit{Flask}\footnote{Ver documentación de Flask : \url{http://flask.pocoo.org/}}.



\subsection{Nanobox}
\textit{Nanobox}\footnote{Documentación de Nanobox : \url{https://docs.nanobox.io/}} es un servicio que emplea Virtual Box y Docker para crear entornos de desarrollo virtuales dentro de la máquina local. Se configura a través de un fichero yaml llamado \textit{boxfile.yml}, en el cual especificaremos la configuración de los recursos que vayamos a necesitar tanto localmente como en producción. 

Tiene una interfaz muy amigable pero la instalación de las dependencias es bastante pesada. Permite desplegar nuevas versiones de forma muy sencilla y transparente para el programador. Tiene muchas funcionalidades de control de memoria, consola para monitorizar los procesos y aunque es de pago, con el paquete de estudiantes se puede probar gratis durante unos meses.

Se descartó porque no me permitía acceder a los permisos de las carpetas que creaba en el servidor, más adelante, hablaremos de ello y de las soluciones que se dieron.

\subsection{Pythonanywhere}
\textit{Pythonanywhere} \footnote{Página principal de Pythonanywhere: \url{https://www.pythonanywhere.com/}} es un servidor web, dedicado exclusiamente a alojar aplicaciones escritas en \textit{Python}. Su interfaz es bastante fácil de entender. Es exclusiva para \textit{Python}, por lo que resulta muy práctico ya que \textit{Flask} es \textit{Python}.

Este hosting nos ofrece dos consolas de \textit{Python} con las que poder interactuar con el proyecto. Se necesita tener el proyecto en un repositorio público y hay que utilizar git para llevarse dicho proyecto al servidor. Una vez clonado el proyecto solo habrá que sincronizarlo y relanzarlo cada vez que queramos un cambio en la aplicación.

\imagen{py}{Interfaz de PYthonanywhere}

\subsection{Comparación entre \textit{Pythonanywhere} y \textit{Nanobox}}
En este apartado se va a hacer una comparativa entre \textit{Pythonanywhere} y \textit{Nanobox} con el fin de argumentar mi decisión final de decantarme por \textit{Pythonanywhere}:

\begin{itemize}
	\item \textbf{Interfaz:} Ambas son bastante intuitivas, pero la de \textit{Nanobox} tiene una interfaz más limpia y bonita a pesar de tener muchas más funcionalidades.
	
	\item \textbf{Precio:} \textit{Nanobox} es de pago y es caro y la preba gratis dura poco. \textit{Pythonanywhere} es gratuito para aplicaciones de poco tránsito.
	
	\item \textbf{Documentación:} \textit{Nanobox} tiene una extensa documentación plagada de ejemplos y tutoriales muy fáciles de seguir. \textit{Pythonanywhere} tiene mucha documentación propia pero su verdadero valor reside en la comunidad que lo usa.
	
	\item \textbf{Información externa a la aplicación:} \textit{Nonobox} no tiene a penas tutoriales externos ni una comunidad que le de un mínimo soporte y por otra parte \textit{Pythonanywhere} tiene una extensa comunidad que comenta dudas, hace tutoriales y por lo general mantienen con vida estas prácticas.
	
	\item \textbf{Funcionamiento:} Con\textit{Pythonanywhere} no tuve a penas problemas, y los que tuve fueron triviales y de fácil solución. Por otra parte \textit{Nanobox} no llegó a funcionarme completamente del todo.
	
\end{itemize}

Mirando los puntos expuestos de forma global, el único que importa es el de \textbf{funcionamiento}, los demás son en mayor o menor medida complementarios.

Por esta razón se ha escogido \textit{Pythonanywhere}.


\section{Lenguaje de programación}
\subsection{Python}
Se ha usado \textit{Python} en su versión más reciente, la 3.7.0 \cite{python}.

Es un lenguaje muy extendido por ser de código abierto, fácil de leer e interpretar.
Es de propósito general, lo que significa que puede abordar muchos frentes en el desarrollo. Es un lenguaje interpretado, multiparadigma (orientación a objetos, programación funcional y programación imperativa), es de tipado dinámico y es independiente de la plataforma en la que se ejecute.

Presenta múltiples ventajas frente a otros lenguajes, entre ellos la gran cantidad de librerías de licencia gratuita que tiene, las funcionalidades que ofrece y la versatilidad de sus funcionalidades.

En este caso es interesante para el proyecto por dos aspectos: 
\begin{itemize}
	\item El despliegue de la aplicación, para el que utilizaremos \textit{Flask} como framework 
	\item \textit{Tensorflow}\footnote{Enlace para el repositorio de Tensorflow: \url{https://github.com/tensorflow}} y lectura de datos. \textit{Python} facilita la lectura de archivos con los que vamos a trabajar en este caso: csv y json, los cuales son usados frecuentemente con \textit{Tensorflow}, la librería que se va a usar para tratar de entrenar un modelo. 
\end{itemize}
\subsubsection{\textit{Anaconda}}
\textit{Anaconda} es una distribución de licencia libre de \textit{Python} y R~\cite{ana}.

Las ventajas del uso de \textit{Anaconda} empiezan desde la gestión de paquetes y librerías para \textit{Python}, gracias a su sistema de administración de paquetes \textbf{Conda} \footnote{Para conocer un poco más conda: \url{https://es.wikipedia.org/wiki/Sistema_de_gesti\%C3\%B3n_de_paquetes}}
lo cual nos facilitará mucho usar librerías enfocadas al procesamiento de datos y en el aprendizaje automático, como Tensorflow, Scikit-team y SciPy.


\subsection{Desarrollo web: \textit{HTML y bootstrap}}
\textbf{HTML} es un lenguaje de marcado, cuyas siglas en inglés son \textit{HyperText Markup Language}~\cite{html}, es decir, lenguaje de marcas de hypertexto.

Se utiliza para confeccionar paginas web. Es un lenguaje basado en la diferenciación, es decir, separar aquellos elementos externos (\textit{Javasecript, CSS}, imágenes, enlaces...) llamándolos a través de referencias a su ubicación.

Con \textit{HTML} se busca que una página escrita en este lenguaje tenga el mismo aspecto y funcionalidad en cualquier buscador en el que se ejecute.

Por otro lado \textbf{Bootstrap}, es un  framework web multiplataforma que contiene plantillas para diseñar aplicaciones más funcionales de forma sencilla~\cite{bootstrap}.

Gracias a estas librerías la página se verá más limpia y profesional sin la necesidad de añadir muchas líneas de \textit{CSS} o \textit{Javascript}.
 

\subsection{Gestión del proyecto}
\subsubsection{Control de versiones, repositorio de GitHub}

Debido a que a lo largo de la carrera nos hemos decantado por el uso de esta herramienta y ya estaba familiarizada con ella será lo que vayamos a utilizar para alojar el proyecto en un repositorio y poder acceder a él de forma controlada por medio de un cliente de \textit{git}.

\subsubsection{ZenHub}
Es una extensión del navegador que se integra con GitHub para poder controlar las tareas que hay que llevar a cabo en el proyecto de forma ágil y rápida de ver.
\imagen{zenhub}{Vista del tablero que nos ofrece la herramienta de ZenHub}

\subsection{Documntación con \LaTeX{}}
\LaTeX{}~\cite{tex} es un lenguaje de marcado mediante el cual escribiremos la memoria y los anexos de la documentación del proyecto.

Es fácil de usar y nos ofrece bastantes ventajas frente a otros editores de texto en cuanto a la personalización. 

Por otro lado la curvatura de aprendizaje de \LaTeX{} puede ser pronunciada al principio.

El entorno en el que se va a editar con \LaTeX{} es \textit{Texstudio}. \footnote{Enlace a la página de Texstudio \url{https://www.texstudio.org/}}

\subsection{Otras librerías}

\subsubsection{Flask Dance}
Para hacer el login se ha usado una librería de Flask que te permite generar un control de acceso mediante otras aplicaciones como \textit{GitHub}, \textit{Twitter} o \textit{Facebook} \cite{dance}.
Se genera un cliente en la API\footnote{La API de Google con la que se ha generado el usuario para la aplicación: \url{https://www.googleapis.com/auth/userinfo.profile}} de Google y se añade para qué direcciones se quiere restringir el acceso, es decir, hay que asociarle las rutas del proyecto.





\capitulo{5}{Aspectos relevantes del desarrollo del proyecto}\label{c.5}

\section{Elección del etiquetador}

Al principio del proyecto se pensó en dividir la funcionalidad del etiquetador de la del clasificador, por lo que una de las primeras tareas fue la de buscar uno que fuese fácil de documentar, fácil de entender por el usuario que vaya a generar las etiquetas, fácil de instalar y que devolviese el conjunto de datos de tal forma que la lectura sea sencilla de hacer desde un script de python.

Después de investigar qe etiquetaro era e más decuado para esta tarea, se optó por una llamada \textit{Alp's Labeling Tools for Deep Learning}\footnote{Enlace para ver la documentación del etiquetador \url{https://alpslabel.wordpress.com/}}
Al final se tomó la decisión de preparar una máquina virtual con el etiquetador instalado y hacer un tutorial por escrito para que el usuario aprendiese a generar las etiquetas como se puede apreciar en la imagen \ref{fig:etiquetapingu}

\imagen{etiquetapingu}{Captura del primer prototipo del etiquetador}

La idea terminó desechándose una vez estuvo finalizada porque la entrada de datos y la nomenclatura de las etiquetas podría ser propensa a errores debido a que depende mucho de que el usuario que etiquete no se equivoque escribiendo el nombre de la etiqueta, además de que la labor llegaba a ser muy pesada.

Esto podría traer problemas con la lectura de datos en el futuro además de pérdida de información.

Llegado a este punto se optó por ofrecer un etiquetador personalizado dentro de la aplicación el cual se puede ver en la imagen \ref{fig:etiquetad}
\imagen{etiquetad}{Vista del etiquetador personalizado.} 

\subsection{Problemas con el etiquetador}
El etiquetador que se proporciona en la aplcación está basado en otro de código abierto llamado \textit{VGG Image Annotator (VIA)}. 

Se escogió porque ofrecía muchas opciones de carga y descarga y modificar la salida de datos para que el \textit{csv} fuese compatible con el clasificador y porque estaba escrito en \textit{JavaScript} completamente. Esto de que estuviese en \textit{JavaScript} era importante porque ha facilitado su integración en el proyecto, no obstante hubo una serie de obstáculos a mencionar:

\begin{enumerate}
	\item En primer lugar, se trata de una aplicación muy grande y con muchas funcionalidades que a priori no se necesitaban, lo cual complicaba la comprensión del código.
	
	\item No había a penas comentarios en todo el código, había muchas funciones que no se utilizaban y los nombres de las variables, en muchas ocasiones, eran poco o nada descriptivas. 
	
	El código del etiquetador es correcto, pero tiene lagunas funcionales que se han tenido que solucionar.
	
	\item A la hora de introducir las etiquetas, el etiquetador original necesita que las introduzcas a mano, pero es precisamente esto de lo que huíamos la última vez, así que se optó por ponerle botones con el nombre ya predefinido de los fitolitos de interés para ahorrarnos erores futuros.
	
	\item Redimensionamiento de las etiquetas. Para recoger los datos en el \textit{csv} existía una problemática porque si cambiabas una etiqueta de tamaño el valor no se actualizaba correctamente el archivo donde se recogían las coordenadas de la etiqueta.
	
	Para solucionar esto se detecto el problema mediante pruebas de etiquetado y se arregló permitiendo que se actualizasen los valores.
	
	\item La falta de experiencia por mi parte en lenguajes como \textit{JavaScript} y \textit{HTML} hizo mella en los tiempos dedicados al etiquetador. Por suerte la curva de aprendizaje es ligera.
	
	\item El diseño de la interfaz. Fue un problema porque al principio hubo problemas con el despliegue, como más adelante explicaremos, y al no poder acceder a los permisos de las carpetas, no podíamos guardar los datos en el servidor, así que para que esto no fuese un impedimento a la hora de obtener las etiquetas, hicimos que no se guardasen en el servidor ni las etiquetas ni las imágenes, pero estas no las descargábamos porque el usuario las subía desde su propio ordenador, por lo que es innecesario.
\end{enumerate}

Al final el etiquetador funciona como se esperaba que lo hiciese.

\section{Problemas en el depliegue}

La primera opción que se usó para alojar la aplicación fue  \textit{Nanobox}. La conforman un conjunto de herramientas que facilitan el mantenimiento de la aplicación. En mi caso generó problemas desde el principio.

\subsection{\textit{Nanobox}: Instalación de \textit{Docker}}
\textit{Nanobox} requiere de un grupo de docker al que se le otorgan privilegios para manipular carpetas. En mi caso no hubo forma de otorgárselos a la primera. Traté de hacerlo por consola pero tampoco los aceptaba.

Borré toda la instalación que había hecha hasta el momento de \textit{Nanobox} y la reinstalé, y conseguí alojar la aplicación, aun así, los permisos de las carpetas en los que se querían guardar las imágenes y etiquetas seguían siendo insuficientes y no hubo forma de cambiarlo.

La siguiente opción fue la dejarlo en local para la presentación, así al menos se verán sus funcionalidades al completo sn depender de permisos en carpetas que no podemos controlar.

En retrospectiva esto fue bueno en varios sentidos:

\begin{itemize}
	\item El primero es que este servidor solo era gratuito porque contaba con un paquete de ayuda para usuarios estudiantes lo que significaba que en octubre caducaba y empezaría a ser de pago sin previo aviso, por lo que en el fondo lo considero una preocupación menos.
	
	\item El Hosting por el que opté después es mucho más cómodo y rápido y aunque tenga menos opciones de mantenimiento y depuración, la prefiero por la simpleza de su mecanismo basado en \textit{Git}.
	
	\item He tenido la oportunidad de probar dos opciones distintas para el despliegue lo cual me ha ayudado a ser más crítica con qué herramientas he de tener en consideración y hsta qué punto es necesario gastar tiempo en ellas.
	
	\item Evito las incompatibilidades del fichero  \textit{yml} y los ficheros de configuración del despliegue, que se quedan en el repositorio. 
\end{itemize}

También ha tenido sus inconvenientes:

\begin{itemize}
	\item Se ha perdido mucho tiempo en algo que sí que aporta valor al proyecto pero no a costa de perder tanto el tiempo en ello. Lo mejor huvise sido haber dejado de lado esta forma de intentar desplegar y haber pasado a la sigiente, para no obcecarme.
	
	\item Se han perdido todas las opciones de mantenimiento, depuración y control de recursos que ofrecía \textit{nanobox}, que no son necesarios a priori, pero no estaba de más cuando había errores o cuando se quería gestionar los recursos que consumía la aplicación.
\end{itemize}

\subsection{Solución: Cambio a \textit{Pythonanywhere}}

Después de mucho tiempo perdido en el despliegue, se decidió que de momento trabajaría en local para poder avanzar. 

Estando más avanzado el proyecto, y con el fin de dejarlo alojado en un servidor, se tomó la decisión de probar otra forma de despliegue.

Pythonanywhere ofreció un despliegue sencillo mediante una consola de bash que hace uso de git.

\section{El clasificador}

Uno de los frentes abiertos del proyecto es la futura integración de un clasificador que permita al usuario subir imágenes que contengan fitolitos y que el sistema de devuelva las imágenes con las etiquetas correspondientes en los fitolitos.

Se ha explorado la opción de hacerlo usando modelos ya hechos proporcionados por \textit{Tensorflow}.

\section{Problemas con las versiones: CUDA, cuDNN y TensorFlow}
Un requisito indispensable es instalar la librería   \textit{TensorFlow-GPU}, debido a que vamos a tratar con imágenes, irá más rápido que con CPU.

Análogamente a esto hay que tener instalado CUDA y cuDNN\footnote{Acceso a la página de nvidia: \url{https://developer.nvidia.com/cudnn}} teniendo cuidado con las versiones tanto de \textit{TensorFlow} como de estas dos anteriores. No es extraño que haya incompatibilidades entre las versiones más recientes. Es mejor informarse bien antes de instalar ninguno de estos elementos.

Por lo general, en la página oficial de \textit{Tensorflow} te especifcan qué versiones son compatibles con qué otros sistemas\cite{ten}.

\section{Entorno Virtual}
Como seguramente haya que estar calibrando las versiones de las librerías que se vayan a instalar, es altamente recomendable crear un entorno virtual para que no haya problemas, ni en el proyecto en el que se está trabajando, ni en los proyectos ajenos a este.

Por lo general, lo ideal es trabajar con versiones actualizadas, pero en este caso es mejor tener la documentación a mano y comprender qué necesidades tiene el clasificador que se quiere entrenar.

\section{Repositorio de TensorFlow}
En las pruebas que se han hecho se ha trabajado con el repositorio de \textit{GitHub} de \textit{TensorFlow}\footnote{Repositorio de TensorFlow: \url{https://github.com/tensorflow/models}}

En él podemos encontrar todos los recursos necesarios para entrenar un modelo.
Se puede clonar o descargar para poder hacer uso de sus funciones.

En este repositorio hay ejemplos de uso para las diferentes funcionalidades que ofrece, por lo que es recomendable saber de antemano qué recursos necesitas para la finalidad que se desea.

\section{Modelos pre-entrenados}
Se ha intentado entrenar al clasificador con dos tipos de modelos\footnote{Enlace para acceder a los modelos: \url{https://github.com/tensorflow/models/blob/master/research/object_detection/g3doc/detection_model_zoo.md}} distintos, pero realmente Tensorflow nos ofrece muchas alternativas, que pueden ser más o meno precisas y cuya velocidad de entrenamiento variará de unas a otras (ver imagen \ref{fig:tabla}).

\imagen{tabla}{Modelos ya entrenados que TensorFlow pone a disposición pública}

\section{Entrenamiento}
Teniendo los conjuntos de imágenes ya etiquetadas y estando generados todos los recursos necesarios para empezar con el entrenamiento, el programa falla debido a una incompatibilidad entre los tipos de datos que se procesan.

Existen varias posibilidades desde mi punto de vista:

\begin{itemize}
	\item Que el problema esté entre las versiones de las librerías de las que se hace uso. 
	
	\item Que el etiquetador esté recogiendo un tipo de dato que no es compatible.
	
	\item Que los modelos con los que se ha probado el entrenamiento procesen un tipo distinto de datos.
	
	\item Que me haya equivocado en cualquier parte del proceso previo al entrenamiento.
	
\end{itemize}

En cualquier caso, el resultado ha sido que no ha podido entrenarse el modelo.

Para poner un ejemplo ilustrativo a cerca de como debería haberse clasificado una imagen, podemos observar el de la imagen \ref{fig:negra}.
Este clasificador ha sido entrenado con el modelo \textit{ssd mobilenet v1 coco}, uno de los más rápidos y que menos potencia necesita, pero también uno de los menos precisos.
\imagen{negra}{Ficha negra clasificada correctamente con el clasificador}

\capitulo{6}{Trabajos relacionados}

Este apartado sería parecido a un estado del arte de una tesis o tesina. En un trabajo final grado no parece obligada su presencia, aunque se puede dejar a juicio del tutor el incluir un pequeño resumen comentado de los trabajos y proyectos ya realizados en el campo del proyecto en curso. 

\capitulo{7}{Conclusiones y Líneas de trabajo futuras}

\section{Conclusiones}
Se han presentado muchos problemas desde el inicio del proyecto en casi todos los aspectos técnicos que he tenido que afrontar.
El primer problema al que me tuve que enfrentar fue al planteamiento de un etiquetador que recogiese de forma segura todo los datos en el formato correcto. Después de probar e investigar al final se decidió personalizar uno y dejarlo funcional.

Otro de los problemas que estuvo presente durante el despliegue de la aplicación con Nanobox, desde la instalación de los requisitos hasta los permisos de acceso a las carpetas una vez estuvo hosteado. Al final la solución consistió en cambiar de método de despliegue.

El mas pesado de los obstáculos ha sido sin ninguna duda el clasificador. Si bien es cierto que he aprendido a reducir un error, a reconocerlo y a diagnosticarlo, me hubiese gustado sacarlo adelante.

Por otro lado he aprendido mucho a cerca de la arquitectura cliente-servidor, he adquirido conocimientos en \textit{HTML} y en \textit{JavaScript}, nuevas librerías de \textit{Python}, y aunque me hubiese gustado aprender más, también he aprendido algo de \textit{Tensorflow}.

A nivel técnico he aprendido bastantes cosas, pero personalmente creo que la más importante ha sido a gestionar mi tiempo y conocer mi ritmo de trabajo.

Ha supuesto todo un reto para mí compaginarlo con el resto del curso, las labores de delegación y ya en verano, con las práctica extracurriculares. 

Aunque sigo sin estar satisfecha con el resultado práctico del trabajo y siendo consciente de que hay muchas cosas mejorables, sí que estoy satisfecha con el hecho de haber puesto a prueba los conocimientos obtenidos a lo largo de la carrera y por supuesto, con todos los conocimientos nuevos que he adquirido.


\section{Líneas de trabajo futuras}

La principal línea de trabajo que veo es la conseguir entrenar al clasificador y ponerlo en producción en el servidor.
Si se sigue usando Tensorflow como lo estaba haciendo yo, ya se dispondría del material necesario para el entrenamiento (en el caso de que el material sea bueno).

Aunque probablemente sea mejor explorar otros tipos de frentes. Existen más librerías a parte de TensorFlow, por ejemplo Keras.

El etiquetador es mejorable, estéticamente no me gusta como escala cuando la imagen que le pasas es demasiado grande.
Sería deseable que se añadiese una opción para poder subir imágenes y etiquetas, a veces el conjunto es muy grande y el usuario puede decidir guardar los cambios y seguir en otro momento.

Quizá sería buena idea plantear un sistema de proyectos para que el usuario pueda guardar diferentes conjuntos de entrenamiento clasificar diferentes conjuntos de imágenes. En ese caso lo más seguro es que hubiese que cambiar de servidor a otro con mayor capacidad.

La galería es una parte de mi aplicación en la que considero que se pueden meter muchas más cosas de carácter funcional, como por ejemplo, eliminar fotos.




\bibliographystyle{plain}
\bibliography{bibliografia}

\end{document}
